% Presentación Final (15 minutos máximo). No es necesario presentar el
% problema asignado. Puntos mínimos exigidos por las reglas:
%   - Representación: variables y movimientos/operadores.
%   - Explicación detallada del código (resaltar correcciones pedidas por las
%     ayudantes en el Estado de Avance).
%   - Aportes de la implementación.
%   - Tablas y/o gráficos comparativos de los resultados.
%   - Conclusiones.
\documentclass[aspectratio=169, 10pt]{beamer}
\usepackage[utf8]{inputenc}
\usepackage[T1]{fontenc}
\usepackage[spanish]{babel}
\usepackage{amsmath}
\usepackage{graphicx}

\usetheme{Madrid}
\setbeamertemplate{bibliography item}[text]

\title{Segmentación multiumbral de imágenes MRI}
\subtitle{Greedy + Tabú Search --- Presentación final}
\author{Matías Peñaloza}
\institute{Inteligencia Artificial 2026-2}
\date{\today}

\begin{document}

\begin{frame}
  \titlepage
\end{frame}

\begin{frame}{Representación y movimientos}
  % TODO: variables, dominio y movimientos/operadores utilizados.
\end{frame}

\begin{frame}{Explicación del código}
  % TODO: solución inicial, función de evaluación (tablas precalculadas,
  %       \cite{Liao2001Fast}), Tabú Search.
\end{frame}

\begin{frame}{Correcciones del Estado de Avance}
  % TODO: una respuesta por cada aspecto indicado por las ayudantes.
\end{frame}

\begin{frame}{Aportes de la implementación}
  % TODO: qué aporta frente al estado del arte~\cite{Hammouche2010Comparative}.
\end{frame}

\begin{frame}{Resultados}
  % TODO: tablas y gráficos comparativos.
\end{frame}

\begin{frame}{Conclusiones}
  % TODO
\end{frame}

\begin{frame}[allowframebreaks]{Referencias}
  \nocite{Otsu1979, Hammouche2010Comparative, Liao2001Fast}
  \bibliographystyle{plain}
  \bibliography{Referencias}
\end{frame}

\end{document}
